% W15 (2026-10-12, restructure P1): Appendix A = Algorithm 1 (old 4), the old Appendix D, and the exactness detail of old 4.2.
\FloatBarrier
\section{PMRT algorithm and proofs}
% W14 (appendix adjudication 2026-10-12, D): retitled from "Proof of Proposition 1 and Corollary 1"; the
% appendix also proves Lemmas 1 to 3.
\label{app:proof}

% Appendix D from the W3 draft (reviews/2026-10-10/drafts/proposition.tex, part (b)), 2026-10-11; Hall and Heyde cited
% at chapter level (RULINGS 4: theorem numbers unverified). Code map: drafts/proposition_notes.md (B).

% W15 (2026-10-12, restructure P1): Algorithm 1 moved from old Section 4.
% W17 (2026-10-12, P2): Algorithm 1 introduced by one sentence and placed before the proofs, so that it no longer floats
% ahead of the appendix text.
Algorithm~\ref{alg:pmrt} collects the steps of Sections~\ref{sec:pmrt} and~\ref{sec:inference} for one hypothesis,
and the proofs below use the notation of Section~\ref{sec:inference}. % src: plan App A (Algorithm 1 here)

\begin{algorithm}
\caption{PMRT for family $f$ and hypothesis $(f, r, k)$}
\label{alg:pmrt}
\begin{algorithmic}[1]
\Require units of $f$ with complete 90~s pre- and post-windows, in information order; propensities $\pi_u$; the fitted parameters (receiver ranking, ridge coefficients $\Gamma$, kernels $\psi_A$ per bin set $A$, clip scale model, training orientation $o$, prior weights $\omega_j$ and directions $d_j$)
\If{fewer than 30 units, 5 accepted, 5 rejected or 3 episodes}
  \State \Return untested (never declared)
\EndIf
\For{each unit $u$}
  \State $v_u \gets \mathrm{sgn}_u\,(L(W_u) - \pi_u)$ \Comment{Equation~\eqref{eq:score}}
  \State form the post-window bins $P_u$ of KPI $k$ per receiver slot (Section~\ref{sec:pmrt}) over the bins $A$ the episode allows (0--90~s, plus 90--150~s where present), and the pre-window features $X_u$
  \State $E_u \gets P_u - \Gamma X_u$ \Comment{ridge prediction, fitted on training episodes}
  \State $g_u \gets \psi_A^{\top} E_u$ with $\psi_A = \Sigma_{AA}^{-1} \mu_A$ \Comment{a kernel per bin set}
  \State $w_u \gets s_u\,\mathrm{clip}\bigl((g_u - \bar g_u) / s_u, \pm c\bigr)$ \Comment{$\bar g_u$ the running mean of earlier scores; Huber clip with a predictable scale $s_u$}
\EndFor
\State $S \gets \sum_u v_u w_u$;\; $\sigma^2 \gets \sum_u \pi_u (1 - \pi_u)\, w_u^2$;\; $z \gets S / \sigma$ \Comment{Equation~\eqref{eq:stat}}
\For{$b = 1, \dots, B$}
  \State redraw $W_u^{(b)} = \text{accept}$ with probability $\pi_u$ for the units of $f$, other families fixed; recompute $v^{(b)}$ and $z_b$ with the same $w$
\EndFor
\State $p \gets$ two-sided value of Equation~\eqref{eq:pval}
\If{$o$ is set}
  \State $p_{+1}, p_{-1} \gets$ one-sided values of Equation~\eqref{eq:pval};\; sign $\gets \mathrm{sign}(S)\, o$
\Else
  \State $p_{+1}, p_{-1} \gets$ unavailable;\; sign $\gets$ sign of the unfiltered statistic
\EndIf
\State \Return $(p, p_{+1}, p_{-1}, \text{sign})$ to the declaration layer (Section~\ref{sec:inference})
\end{algorithmic}
\end{algorithm}

Table~\ref{tab:pipeline} lists the construction of the weight $w_u$ of unit $u$ for hypothesis $(f, r, k)$ (steps 1 to 7, in order) and the inputs of the declaration layer (step 8). Training used 480 earlier and 20 development episodes. Step 8 runs steps 1 to 7 on the 480 earlier episodes and uses their $z$ of Equation~\eqref{eq:stat}; $z_{\mathrm{train}}$ is the $z$ of the plain 0--90~s total of the step-3 residuals on the training episodes. Every fitted part was fixed before any evaluation data were generated.

\begin{table}[tbp]
\caption{Weight construction.}
\label{tab:pipeline}
% W17 P2b: cells shortened, the "Frozen" column dropped (it repeated the caption; step 5 says "online" in "Fitted on").
% Dropped detail, kept in the sources below and in Algorithm 1: "a unit needs the first 90 s" (Section 4.2 window),
% "projected to the number of analysed episodes" (FL prior_weights).
% Row 8 (W17, moved from the old 4.2 prose): FL prior_weights (Roeder-Wasserman, floor .2, cap W_CAP 5, mean 1, projected
% to the analysed episodes); FL prior_directions (|prior z| >= 3); PV4 s.3 (33 of 60); FL docstring (independent data).
% Hand-written from the old Section 4.1 prose and its footnote (W17, plan 4.3; new-reviewer-2 s.4e). Sources per row:
% PM docstring (POST bins [0,30) [30,60) [60,90) [90,150) s after t0; slots: fixed receiver score rho_f from f's effect
% on per-cell load, top-(S-1) cells + the rest, NSLOT own 1, nbr 4, far 2; unit set = crt_units H = 90 set, a unit
% without the [90,150) bin uses the kernel restricted to the observed bins; adjust: X_u = pre bins, other KPIs' pre,
% own-cell pre, ctx, hist; kernel loadsp K = Sigma^-1 mu, Ledoit-Wolf, sign-constrained spatial profile x time profile;
% score r_u = k_u - m_u, running centre, PREDICTABLE; weights "_c": Huber clip, predicted sd; sign rule |z_plain_train|
% >= 2); PV4 s.3 (artifact holds the kernels with / without the [90, 150) bin; frozen); training: ev2 sub v2 480 eps +
% DEV 20 (PM docstring; TERM s.6.1).
\footnotesize
\setlength{\tabcolsep}{3pt}
\begin{tabular*}{\textwidth}{@{\extracolsep{\fill}}>{\raggedright}p{0.2\textwidth}>{\raggedright}p{0.6\textwidth}>{\raggedright\arraybackslash}p{0.15\textwidth}@{}}
\toprule
Step & What it does & Fitted on \\
\midrule
1 Target bins & KPI $k$ over the cells of $r$, at 0--30, 30--60, 60--90 and 90--150~s after opening & -- \\
2 Receiver slots & Own 1, neighbour 4, far 2 slots: cells whose load the family moved most get one each, the rest share one & training \\
3 Adjustment & Ridge prediction~\cite{hoerl1970ridge} from pre-window bins, context and earlier units' centred modes; residual kept & training \\
4 Matched filter & Score $g_u$: step-3 residuals weighted by $\psi = \Sigma^{-1} \mu$, favouring bins with more effect $\mu$ (one sign over slots) and less noise $\Sigma$ (Ledoit--Wolf~\cite{ledoit2004covariance}); restricted to the bins present & training \\
5 Score centre & Subtracts $\bar g_u$, the running mean score of earlier units & earlier units, online \\
6 Huber clip & Clips at a scale predicted at opening~\cite{huber1964robust}, giving $w_u$ & training \\
7 Orientation & $o$: sign of $\mu$ summed over 0--90~s, set only if $|z_{\mathrm{train}}| \ge 2$ & training \\
\midrule
8 Declaration weights & Roeder--Wasserman $\omega_j$~\cite{roeder2009weighted}: floor 0.2, cap 5, mean 1; direction $d_j$, sign of $z$ where $|z| \ge 3$ (33 of 60) & the 480 earlier episodes \\
\botrule
\end{tabular*}
\end{table}


\subsection{Setting}
\label{app:proof_setting}

% W17 P2b (ruling Q1): the sharp-null and G_u paragraph of Section 4.4 moved here (first two sentences and the
% last sentence of this paragraph). src: W3 draft (null); proposition_notes.md C1, C2; V2 B22 (invariance, not
% pre = post); V2 B4 (statement valid under H0 only)
Fix a hypothesis $(f, r, k)$. For every key $\xi$ defined below, let $Y_\xi$ be the target that a unit of $f$ with
that key would have: KPI $k$ over the cells of relation $r$ in the pre- and post-windows around its opening time. The
sharp null $H_0(f, r, k)$ states that every $Y_\xi$ is the same under any assignment of $f$'s modes. It is stated over
keys rather than realised units, because which units are realised depends on the modes. In
Section~\ref{sec:inference}, $\mathcal{G}_0$ holds all randomness of the system except the draws that set the modes of
$f$, and $\mathcal{G}_u$ adds the modes $W_1, \dots, W_u$. The probability space carries the network's randomness and the
assignment draws of every family. The draws of family $f$ form a field $\{U_\xi\}$ of
independent Uniform$(0, 1)$ variables indexed by keys $\xi$ (episode, cell, xApp, opening
second), independent of $\mathcal{G}_0$, which holds all other randomness. Unit $u$ of
$f$, in information order (episode, opening time, cell), has key $\xi_u$ and mode
$W_u = \text{accept}$ if and only if $U_{\xi_u} < \pi_u$, and
$\mathcal{G}_u = \mathcal{G}_0 \vee \sigma(U_{\xi_1}, \dots, U_{\xi_u})$. This filtration
contains the one of Section~\ref{sec:inference}, built from the modes; every statement below
holds for both, because keys, propensities and weights depend on the draws only through the
modes. Distinct units of
$f$ have distinct keys, because a cell-xApp pair has at most one active unit. Episodes are
appended in a fixed order as $n$ grows, so the units of the first $n$ episodes are an initial
segment of one sequence and $\mathcal{G}_u$ does not depend on $n$. Write
$L_u = L(W_u)$, so that
$v_u = \mathrm{sgn}_u (L_u - \pi_u)$, with conditional variance $\pi_u (1 - \pi_u)$. Under $H_0(f, r, k)$ every $Y_\xi$ is
$\mathcal{G}_0$-measurable, so the target $Y_{\xi_u}$ of unit $u$, and hence the post-window part of $w_u$, is
$\mathcal{G}_{u-1}$-measurable when $\xi_u$ is, although it is observed after the unit opens. Only then may $w_u$ use
$u$'s own post-window, and validity concerns that case alone.

\begin{lemma}[Fresh draws]
\label{lem:fresh}\claim{THM-FRESH}
If $\xi_u$ and $\pi_u$ are $\mathcal{G}_{u-1}$-measurable and $\xi_u$ differs from
$\xi_1, \dots, \xi_{u-1}$, then $U_{\xi_u}$ is Uniform$(0, 1)$ and independent of
$\mathcal{G}_{u-1}$, so $\Pr(W_u = \text{accept} \mid \mathcal{G}_{u-1}) = \pi_u$.
\end{lemma}

\begin{proof}
Fix a key $\xi$ and let $\mathcal{H}_\xi = \mathcal{G}_0 \vee \sigma(U_{\xi'} : \xi' \ne \xi)$,
which is independent of $U_\xi$. The keys $\xi_1, \dots, \xi_{u-1}$ and the values
$U_{\xi_1}, \dots, U_{\xi_{u-1}}$ are built recursively, each key from $\mathcal{G}_0$ and the
earlier values. On the event $E_\xi = \{\xi_u = \xi\}$ none of the earlier keys equals
$\xi$, so the recursion never reads $U_\xi$, and by induction on $u$ every event
$A \in \mathcal{G}_{u-1}$ satisfies $A \cap E_\xi \in \mathcal{H}_\xi$. Hence, for every Borel
set $C$,
$\Pr(A \cap E_\xi \cap \{U_{\xi_u} \in C\}) = \Pr(A \cap E_\xi)\, \lambda(C)$ with
$\lambda$ the Lebesgue measure, and summing over the countably many keys gives
$\Pr(A \cap \{U_{\xi_u} \in C\}) = \Pr(A)\, \lambda(C)$.
\end{proof}

In the simulator the assumptions of Lemma~\ref{lem:fresh} hold because the trajectory up to an
opening time is a function of $\mathcal{G}_0$ and of the modes of earlier units. A configuration change takes effect at the next tick, the xApps and the incumbent read delivered reports only, and
a unit of $f$ that opens in the same second at a lower-indexed cell precedes $u$ in information
order. Assumption (A1) states this property for a general system. The propensity is 0.5 under the randomised
policy and within $[0.15, 0.85]$ under the incumbent. % src: collect_p.py:95, 118-130, 193-219; units_p.py:236-247 (moved from Section 4.4, W17 P2b)
Assumption (A3) holds by the construction of Table~\ref{tab:pipeline}, and (A4) is the one assumption that the
construction leaves open. % src: W3 draft (A3, A4 status); pmrt.py:173-202, 295-323, 354-381, 581-590 (moved from Section 4.4, W17 P2b)
 The number $N$ of units of $f$ is
bounded by a deterministic $n N_{\max}$, where $N_{\max}$ is the number of cell-xApp pairs times the number of
minutes in an episode, since a pair opens at most one unit per 60~s. Setting $v_u w_u = 0$ for the indices
$N < u \le n N_{\max}$ keeps every statement below, because the event
$\{u \le N\}$ is $\mathcal{G}_{u-1}$-measurable.

\subsection{Assumptions in detail}
\label{app:assumptions}

Assumption (A4) of Section~\ref{sec:inference} states formally that, as $n$, the number of episodes, grows, there are constants $s_n \to \infty$ with $\sigma^2 / s_n^2 \to \eta^2$ in probability and
$0 < \eta^2 < \infty$ almost surely, that $\max_u |w_u| / s_n \to 0$ in probability, and that each $w_u$ is square
integrable; Section~\ref{sec:inference} states what the four assumptions require in operational terms.

\subsection{Proof of Proposition~\ref{prop:pmrt}}
\label{app:proof_prop}\claim{CIT-HH}\claim{CIT-DUR}

\emph{Martingale.} By (A1), $\mathbb{E}[L_u \mid \mathcal{G}_{u-1}] = \pi_u$, so
$\mathbb{E}[v_u \mid \mathcal{G}_{u-1}] = 0$ and
$\mathbb{E}[v_u^2 \mid \mathcal{G}_{u-1}] = \pi_u (1 - \pi_u)$, using
$\mathrm{sgn}_u^2 = 1$. By (A3), $w_u$ is $\mathcal{G}_{u-1}$-measurable, so
$\zeta_u = v_u w_u$ satisfies $\mathbb{E}[\zeta_u \mid \mathcal{G}_{u-1}] = 0$ and
$\sum_u \mathbb{E}[\zeta_u^2 \mid \mathcal{G}_{u-1}] = \sum_u \pi_u (1 - \pi_u)\, w_u^2 = \sigma^2$.
These identities are exact for every $n$.

\emph{Part 1.} Let $X_{n,u} = \zeta_u / s_n$. Since $|v_u| \le 1$,
\begin{equation*}
\sum_u \mathbb{E}\bigl[X_{n,u}^2 \mathbf{1}\{|X_{n,u}| > \gamma\} \mid \mathcal{G}_{u-1}\bigr]
\le \frac{\sigma^2}{s_n^2}\, \mathbf{1}\Bigl\{\max_u |w_u| > \gamma s_n\Bigr\} \to 0
\end{equation*}
in probability for every $\gamma > 0$ by (A4), and
$\sum_u \mathbb{E}[X_{n,u}^2 \mid \mathcal{G}_{u-1}] = \sigma^2 / s_n^2 \to \eta^2$ in probability.
The histories $\mathcal{G}_u$ do not depend on $n$, so they are nested in $n$, and the rows have the
deterministic length $n N_{\max}$ (Appendix~\ref{app:proof_setting}). The martingale central limit theorem with a random
limiting variance~\cite[Cor.~3.1]{hall1980martingale} then gives $S / s_n \to \eta Z$ stably, with
$Z \sim \mathcal{N}(0, 1)$ independent of $\eta$, and since $\eta^2 > 0$ almost surely its self-normalised
form~\cite[Cor.~3.2]{hall1980martingale} gives $z = S / \sigma \to Z$ in distribution. The event $\sigma = 0$, on which the code
sets $z = 0$ and $p = 1$, has probability tending to 0.

\emph{Part 2.} Condition on the data. A redraw sets $W_u^{(b)} = \text{accept}$ with probability
$\pi_u$, independently over $u$ and of the data, and keeps $\mathrm{sgn}_u$, $\pi_u$, $w_u$ and
$\sigma$. Thus $z_b = \sigma^{-1} \sum_u w_u v_u^{(b)}$ is a sum of independent terms with
mean 0 and variances $w_u^2 \pi_u (1 - \pi_u) / \sigma^2$ summing to 1. Its Lindeberg
quantity is at most $\mathbf{1}\{\max_u |w_u| > \gamma \sigma\}$, which tends to 0 in
probability because $\max_u |w_u| / s_n \to 0$ and $\sigma / s_n \to \eta > 0$. Along a further subsequence of every subsequence it tends to 0 almost surely for every
rational $\gamma$, and there the Lindeberg--Feller theorem~\cite[Thm.~3.4.10]{durrett2019probability} and
P\'olya's theorem~\cite[Exercise~3.2.9]{durrett2019probability} give $\sup_x |\Pr^*(z_b \le x) - \Phi(x)| \to 0$
almost surely, where $\Pr^*$ is the redraw law; hence the convergence holds in probability.

\emph{Part 3.} The code counts $|z_b| \ge |z| - 10^{-9}(1 + |z|)$; the tolerance
can only increase $p$, so it suffices to prove the claim for the count $Q$ without it, which
is Equation~\eqref{eq:pval}. Given the data, the $B$ redraws are independent and identically
distributed, so $Q$ is Binomial$(B, \Pi_n)$ with $\Pi_n = \Pr^*(|z_b| \ge |z|)$. By part 2,
$|\Pi_n - 2\{1 - \Phi(|z|)\}| \to 0$ in probability, and by part 1,
$2\{1 - \Phi(|z|)\}$ converges in distribution to a Uniform$(0, 1)$ variable $V$; so $\Pi_n \to V$ in
distribution. For each $j \in \{0, \dots, B\}$ the map
$x \mapsto \binom{B}{j} x^j (1 - x)^{B - j}$ is bounded and continuous on $[0, 1]$, so
\begin{equation*}
\Pr(Q = j) = \mathbb{E}\Bigl[\binom{B}{j} \Pi_n^j (1 - \Pi_n)^{B - j}\Bigr]
\to \int_0^1 \binom{B}{j} x^j (1 - x)^{B - j}\, dx = \frac{1}{B + 1}.
\end{equation*}
Thus $p = (1 + Q) / (B + 1)$ converges in distribution to the uniform law on
$\{1/(B+1), \dots, 1\}$. Because $p$ takes values in this fixed finite set,
$\Pr(p \le t) \to \lfloor t (B + 1) \rfloor / (B + 1) \le t$ uniformly in $t$. For $p_d$ replace
$\Pi_n$ by $\Pr^*(d\, o\, z_b \ge d\, o\, z)$, which is within $o_p(1)$ of $1 - \Phi(d\, o\, z)$;
since $|d\, o| = 1$, $d\, o\, z \to \mathcal{N}(0, 1)$ and the same argument applies. The last
sentence of the proposition follows from $|\Pi_n - 2\{1 - \Phi(|z|)\}| \to 0$ in probability and
$\mathbb{E}[Q / B \mid \text{data}] = \Pi_n$.
\hfill$\square$

\subsection{When the redraw p-value is exact}
\label{app:proof_exact}

\begin{lemma}[Exact case]
\label{lem:exact}\claim{THM-EXACT}\claim{CIT-PS}
Suppose, in addition to (A1) and (A2), that the realised units, their order, $\mathrm{sgn}_u$,
$\pi_u$ and $w_u$ are all $\mathcal{G}_0$-measurable. Then $\Pr(p \le t \mid \mathcal{G}_0) \le t$
for every $t \in [0, 1]$, every $n$ and every $B$, and likewise for $p_d$.
\end{lemma}

\begin{proof}
Applying Lemma~\ref{lem:fresh} unit by unit, conditionally on $\mathcal{G}_0$ the modes are
independent with propensities $\pi_u$. The observed assignment and the $B$ redraws are then independent
and identically distributed given $\mathcal{G}_0$, and $z$ is the same fixed function of each, so
the $B + 1$ values $|z|, |z_1|, \dots, |z_B|$ are exchangeable given $\mathcal{G}_0$. The Monte Carlo
p-value of Equation~\eqref{eq:pval}, which counts ties against $|z|$, is then valid for every
$B$~\cite[Sect.~4]{phipson2010permutation}.
\end{proof}

PMRT does not meet the extra condition of Lemma~\ref{lem:exact}, even under the sharp null, for three reasons.
(i)~Which units exist and when depends on earlier modes. A rejected request opens a new unit after the hold, and an
accepted change starts an interval in which the actuator refuses a new change and no unit opens. (ii)~Under the
incumbent, $\pi_u$ depends on the delivered state, which earlier modes move. (iii)~$w_u$ contains the design-centred modes of earlier units of $f$, and the
pre-windows of the other KPIs and the context, which earlier modes move unless those KPIs are null
too, and the score centre and the clip scale inherit both. On placebo logs, where accept is
applied to every request, (i) and (ii) disappear, but (iii) remains through the earlier modes in
$w_u$.

% W15 (2026-10-12, restructure P1): exactness detail moved from old 4.2 (L905-913).
% W17: DUP resolved (W15 marker): the moved paragraph repeated the three reasons above; its two remaining points (no
% study is exact; the example) kept as the two sentences below. Section 4.4 keeps the two-sentence summary.
The redraw p-value is therefore exact in none of the datasets, placebo logs included. % src: W3_DONE; RULINGS 5; proposition_notes.md C15 (no exactness claim)
Example~\ref{ex:adaptive} shows a two-unit design in which re-proposal alone makes the one-sided redraw p-value
anti-conservative. % src: W3 draft Example D.1; RULINGS 5

\begin{example}[Re-proposal alone breaks finite-sample validity]
\label{ex:adaptive}\claim{THM-COUNTEREX}
Let $\pi_u = 0.5$, $\mathrm{sgn}_u = +1$ and $o = +1$. Unit 1 has weight $w_1 = 1$. If it is
rejected, the request is proposed again and opens unit 2 with weight $w_2 = a > 1$; if it is
accepted, there is no unit 2. All weights are fixed in advance, so (A1) to (A3) hold. If unit 1 is
accepted, $z = 1$ and the redraws give $z_b = \pm 1$, so $p_{+1} = 1/2$. If unit 1 is rejected,
$z = (a - 1)/\sqrt{1 + a^2}$ when unit 2 is accepted and $z = -(a + 1)/\sqrt{1 + a^2}$ when it is
rejected; the redraws on the two realised units give $p_{+1} = 1/2$ and $p_{+1} = 1$. Hence, with
$B \to \infty$, $\Pr(p_{+1} \le 1/2) = 1/2 + 1/4 = 3/4 > 1/2$. The redraws treat unit 2 as if it
existed whatever unit 1 received, which is the source of the error, as in randomisation tests under
response-adaptive designs~\cite{simon2011randomization}; as the number of units grows,
the martingale central limit theorem removes it.
\end{example}

\subsection{Proof of Corollary~\ref{cor:fdr}}
\label{app:proof_cor}\claim{THM-BY-FIN}\claim{CIT-BR}

Let $c_m = \sum_{i=1}^m 1/i$ and $t_{j,i} = q\, \omega_j\, i / (m c_m)$ for $i = 1, \dots, m$, with $t_{j,0} = 0$. The
declaration layer is the weighted step-up with the reshaping function $i \mapsto i / c_m$, which rejects exactly the hypotheses
with $p_j \le t_{j,\rho}$, $\rho$ being the number of rejections, and which controls the FDR at
$q \sum_{j \in \mathcal{H}_0} \omega_j / m$ under arbitrary dependence when the null p-values are
super-uniform~\cite[Thm.~4.2]{blanchard2008fdr}. The p-value used for $j$, $p_{d_j}$ or $p$, is fixed in advance, and
an untested hypothesis is never rejected, which only lowers the FDR. Here the null p-values are super-uniform only in
the limit, so the bound is re-derived. Since $p_j \le t_{j,\rho}$ with $t_{j,i-1} < p_j \le t_{j,i}$ implies
$\rho \ge i$, the argument of~\cite{benjamini2001dependency} gives, whatever the joint law of the p-values and with
$F_j(t) = \Pr(p_j \le t)$,
\begin{equation*}
\mathbb{E}\Bigl[\frac{\mathbf{1}\{p_j \le t_{j,\rho}\}}{\max(\rho, 1)}\Bigr]
\le \sum_{i=1}^{m-1} F_j(t_{j,i}) \Bigl(\frac{1}{i} - \frac{1}{i+1}\Bigr) + \frac{F_j(t_{j,m})}{m}.
\end{equation*}
For $j \in \mathcal{H}_0$, Proposition~\ref{prop:pmrt} gives
$\limsup_n \sup_t \{F_j(t) - t\} \le 0$, uniformly in $t$, so the grid may move with $n$, as it does
through the projected weights. Substituting $F_j(t_{j,i}) \le t_{j,i} + o(1)$,
\begin{equation*}
\limsup_{n} \mathbb{E}\Bigl[\frac{\mathbf{1}\{p_j \le t_{j,\rho}\}}{\max(\rho, 1)}\Bigr]
\le \frac{q\, \omega_j}{m c_m} \Bigl( \sum_{i=1}^{m-1} \frac{1}{i+1} + 1 \Bigr) = \frac{q\, \omega_j}{m},
\end{equation*}
and summing over $j \in \mathcal{H}_0$ gives Equation~\eqref{eq:fdr}. The weights and directions
enter only through the fixed thresholds $t_{j,i}$, which is why they must not depend on the
evaluation data. \hfill$\square$

\subsection{The design slope}
\label{app:proof_slope}

\begin{lemma}[Design slope]
\label{lem:slope}\claim{THM-SLOPE}
Assume (A1) with the natural history $\mathcal{F}_{u-1}$ in place of $\mathcal{G}_{u-1}$, bounded
outcomes, a centre $\tilde y_u$ that is $\mathcal{F}_{u-1}$-measurable, $N / n$ bounded away from 0 and
infinity with probability tending to one. Then
$\hat\beta$ of Equation~\eqref{eq:slope_ht} satisfies
\begin{equation*}
\hat\beta - \frac{\sum_u \pi_u (1 - \pi_u)\, \mathrm{sgn}_u\, \tau_u^{\pi}}{\sum_u \pi_u (1 - \pi_u)} \to 0
\end{equation*}
in probability. When $\pi_u = 0.5$ for every unit the right-hand fraction is $\bar\tau_{\pi}$
and Equation~\eqref{eq:slope_ht} holds as an identity.
\end{lemma}

\begin{proof}
Since $\tilde y_u$ is $\mathcal{F}_{u-1}$-measurable and $\mathbb{E}[v_u \mid \mathcal{F}_{u-1}] = 0$,
$\mathbb{E}[v_u e_u \mid \mathcal{F}_{u-1}] = \mathrm{sgn}_u \{(1 - \pi_u)\, \mathbb{E}[y_u \mathbf{1}\{W_u = \text{accept}\} \mid \mathcal{F}_{u-1}]
- \pi_u\, \mathbb{E}[y_u \mathbf{1}\{W_u = \text{reject}\} \mid \mathcal{F}_{u-1}]\}
= \pi_u (1 - \pi_u)\, \mathrm{sgn}_u\, \tau_u^{\pi}$, and
$\mathbb{E}[v_u^2 \mid \mathcal{F}_{u-1}] = \pi_u (1 - \pi_u)$. The sums
$\sum_u \{v_u e_u - \pi_u (1 - \pi_u) \mathrm{sgn}_u \tau_u^{\pi}\}$ and
$\sum_u \{v_u^2 - \pi_u (1 - \pi_u)\}$ are martingales with bounded increments over at most $n N_{\max}$ terms (Appendix~\ref{app:proof_setting}), so
their variances are $O(n)$ and, divided by $N$, they tend to 0 in probability by Chebyshev's inequality. Overlap keeps $N^{-1} \sum_u \pi_u (1 - \pi_u) \ge \pi_{\min}(1 - \pi_{\min})$, and bounded outcomes keep
$N^{-1} \sum_u \pi_u (1 - \pi_u) \mathrm{sgn}_u \tau_u^{\pi}$ bounded, which gives the limit of the ratio. With $\pi_u = 0.5$, $v_u^2 = 1/4$ and
$4 v_u = 2\, \mathrm{sgn}_u (2 L_u - 1)$, which is Equation~\eqref{eq:slope_ht}.
\end{proof}

% src: drafts/proposition_notes.md C14 (e6p_conf_analyze.py:201-215; centre eprocess_units.py:132-148); RULINGS 17:
% one-sentence disclosure here only, effect not quantified.
% W17 (2026-10-12, P2): "the effect of this was not quantified" replaced by the X16 result (P2_COMMON "The X16
% result": it replaces "not measured" wherever it appears). src: x16_centre.json (maps rel_change -0.0286 to +0.0647,
% all sign_agree_strict; certification status_changes []); X16_DONE.md.
